% WHAT AN ALIGNMENT TEMPLATE ENDS WITH.
%
% The second half of a column's template ends with a TOKEN of its own --
% a name no document can write, which the reference calls \endtemplate --
% and that token is part of the same list as the rest of the template. Two
% things follow:
%
%   the context names it, so a fault drawn while the entry is finishing
%   reads `<template> \endtemplate' where the column's second half is
%   empty, and `<template> &\ifnum 0=`{\fi \endtemplate' where it is not;
%
%   a macro called at the very end of an entry can take it as its
%   ARGUMENT. `\def\trap#1{}\halign{#\cr\trap\cr}' hands \trap the
%   \endtemplate and draws nothing at all, where an entry that simply ran
%   out would have drawn a runaway.
%
% What the token stands for is the end of the entry's material. It is not
% traced for itself: what the reference traces is the \endv it turns into,
% which is written `end of alignment template'.
%
% A lookahead sees it as any other token -- `\global\futurelet\x\relax&'
% leaves \x meaning `\outer endtemplate:'.
%
% HSTeX ended a template where its tokens ran out, so there was no token to
% name, to hand to a macro, or to look at. trip line 422 writes
% `\e\trap\cr', which is \expandafter\trap\cr at the end of an entry.
\catcode`\{=1 \catcode`\}=2 \catcode`&=4 \catcode`\#=6
\tracingonline=1 \tracingcommands=2 \tracingmacros=2
\hsize=100pt
\message{[A a macro that takes the ending token as its argument]}
\def\trap#1{}
\setbox0=\vbox{\halign{#\cr\trap\cr}}
\message{[B a lookahead at the end of an entry]}
\setbox0=\vbox{\halign{#\cr A\global\futurelet\x\relax\cr}}
\show\x
\message{[C what the context names while an entry finishes]}
\setbox0=\vbox{\halign{#&\undefinedhere#\cr A&B\cr}}
\message{[done]}
\end
